\section{Displacement, measured}
\label{sec:displacement}

This section asks where the loss lands when wages are displaced, at the one displacement level
that sits inside the data used to calibrate reemployment. The answer is that displacement is a
fiscal event well before it is a financial one, and that the composition of the private part of it
depends on how the automation happens rather than only on how much of it there is.

\subsection{The boundary, stated before any number}

The reemployment rate in Section~\ref{sec:data} is a straight line fitted to fourteen biennial
observations and then solved as a fixed point. A fitted line extrapolates until it breaks, and
this one breaks: past a certain displacement level the fixed point leaves the range of
reemployment rates ever observed, and past a further point no admissible solution exists at all.
We report a point estimate only where the solution lies inside the observed range and a band
otherwise, labelled as outside the data.

The physically exposed group leaves the observed range first, above the ten percent level. The two
cognitively exposed groups stay inside it to twenty-five percent. Above about a quarter of the
wage bill displaced, no group has a point estimate, and every figure in this paper at those levels
is a band.

This is a limit of our specification, not a fact about the economy, and it is stated that way
wherever it appears. An earlier version of the code enforced the bound accidentally, through an
iteration that clipped at zero, and published the clipping boundary as though it were an estimate.
Correcting it removed the point estimate from about a quarter of the cells in the response surface.
Where the old value differed from the new band it was the boundary itself, so the superseded
figures overstated the loss at large displacement rather than understating it.

\subsection{First-round incidence by holder}

Table~\ref{tab:incidence} gives the first-round loss by holder at ten percent of the wage bill,
for each of the three exposure groups. The terminal-year fiscal loss at that level is
\result{TerminalLossTen} billion dollars, which is replicated: an independent rebuild landed
within a fifth of a percent of it.

Read across the table, the shape is the same whichever group is displaced. The federal government
takes between \result{FedFirstRoundTenLow} and \result{FedFirstRoundTenHigh} billion dollars of the
first round, private holders between \result{PrivateFirstRoundTenLow} and
\result{PrivateFirstRoundTenHigh} billion. Most of the federal loss is revenue rather than credit:
general revenue, then the social security trust funds at \result{OASDILossTenLow} to
\result{OASDILossTenHigh} billion and the hospital insurance fund at \result{HILossTenLow} to
\result{HILossTenHigh} billion. The credit losses that reach the federal balance sheet, through the
student loan book, the housing administration and the housing agencies, are a minority of the
total.

\begin{table}[htbp]
\centering
\caption{First-round loss by holder at ten percent of the wage bill displaced, billions of
dollars. Measured; the terminal-year fiscal loss is replicated.}
\label{tab:incidence}
\setlength{\tabcolsep}{4pt}
\begin{tabular}{p{4.8cm}p{2.2cm}p{2.5cm}p{2.5cm}}
\toprule
Where the first-round loss lands & Physically exposed & Cognitively exposed, first index & Cognitively exposed, second index \\
\midrule
Federal general revenue & 98.7 & 90.0 & 112.4 \\
Social security trust funds & 63.7 & 54.0 & 72.5 \\
Hospital insurance trust fund & 17.1 & 17.4 & 20.9 \\
Federal student loan book & 5.8 & 12.1 & 11.5 \\
Federal housing administration & 3.2 & 5.7 & 5.9 \\
Housing agencies, gross & 12.9 & 23.2 & 23.9 \\
Federal, total & 188.4 & 179.2 & 223.2 \\
Private holders, total & 32.0 & 49.0 & 52.3 \\
Federal share & 0.855 & 0.785 & 0.810 \\
\bottomrule
\end{tabular}

\end{table}

\subsection{The federal share, by displacement level and by classification}

The federal share of the first round is not a single number, and reporting it as one hides the
one thing it does, which is rise with the displacement level. It also depends on a classification
choice about the housing agencies, and both readings travel together throughout.

Under the \emph{narrow} reading, only the agency loss beyond the agencies' own capital and a year
of earnings counts as federal. Under the \emph{conservatorship} reading, the retained agency loss
counts as federal, because the agencies are in conservatorship, their net worth stands behind
agreements with the Treasury, and no private shareholder is positioned to absorb anything; only
what private credit cover actually pays is private. The narrow reading is the lower bound and the
conservatorship reading the upper bound.

At the ten percent level the federal share is \result{FedShareTenNarrowLow} to
\result{FedShareTenNarrowHigh} percent under the narrow reading and \result{FedShareTenConsLow} to
\result{FedShareTenConsHigh} percent under the conservatorship reading. Table~\ref{tab:fedshare}
gives the whole path. It rises with displacement to a peak around the fifty percent level and then
falls back, because the fiscal component saturates while credit losses keep growing. Those upper
rows are outside the data and are reported for shape rather than for level. The measured direction,
the rise from the five to the twenty-five percent level, was reproduced independently: the
rebuilder's own by-level result sits inside our range at both ends.

\begin{table}[htbp]
\centering
\caption{Federal share of the first-round loss, by displacement level and by the classification of
the housing agencies. Measured inside the data, scenario above it.}
\label{tab:fedshare}
\begin{tabular}{lrrl}
\toprule
Displacement, share of the wage bill & Narrow reading & Conservatorship reading & Inside the data \\
\midrule
5 per cent & 0.759 to 0.853 & 0.857 to 0.901 & yes \\
10 per cent & 0.785 to 0.870 & 0.876 to 0.912 & yes \\
25 per cent & 0.810 to 0.923 & 0.878 to 0.948 & yes \\
50 per cent & 0.855 to 0.925 & 0.897 to 0.949 & no \\
75 per cent & 0.809 to 0.897 & 0.859 to 0.924 & no \\
\bottomrule
\end{tabular}

\end{table}

\subsection{The housing agencies absorb their share without reaching the Treasury}

The agency book was rebuilt loan class by loan class from the two Enterprises' 2025 annual filings,
which is what makes the previous paragraph's classification question answerable rather than
assumed. Three measured results come out of it.

First, a large part of each single-family book carries no credit enhancement at all:
\result{UnenhancedLow} to \result{UnenhancedHigh} percent across the two Enterprises. Mortgage
insurance and credit risk transfer cover the rest, and they cover it thinly.

Second, private cover transfers only \result{PrivateCoverLow} to \result{PrivateCoverHigh} percent
of an agency loss at the ten percent level. The structures are attachment-point structures written
for idiosyncratic losses; a displacement event of this size passes through most of the loss to the
guarantor.

Third, and this is the result that matters for the classification question, the Treasury draw is
\result{TreasuryDrawBn} at every displacement level we ran, including those outside the data:
\result{AgencyCapitalBn} billion dollars of Enterprise capital stands against a maximum retained
loss of \result{MaxRetainedAgencyLossBn} billion. The agencies absorb what reaches them. That is a
first-round result with house prices held fixed, and it is the honest scope of it: a house price
fall would change the loss given default on every loan in the book, and the engine does not model
one.

\subsection{Trust funds}

Trust fund effects are computed on a payroll-only basis. Each fund's loss is the displaced share
of its own covered payroll, adjusted for the wages that come back, times that fund's own payroll
share of its own income. The two funds have different payroll shares and different caps: the social
security funds are capped at the contribution and benefit base, so displacing a high-paid worker
costs them less than displacing a low-paid one, while the hospital insurance fund has no cap. That
asymmetry is why the two rows in Table~\ref{tab:incidence} move differently across the three
exposure groups.

One consequence of the payroll basis is worth carrying into Section~\ref{sec:institutions}:
replacement income funded from a capital tax cannot repair these funds, because it is not covered
wages. A transfer that makes a displaced household whole leaves the trust fund shortfall exactly
where it was.

\subsection{Households, and where the buffers are thinnest}

Table~\ref{tab:buffers} gives the liquid runway of working households by wage quintile: the number
of months of income their liquid assets cover, and the share with less than one month of it.

The finding is that buffers are not monotonic in pay. The thinnest buffers are in the
\emph{second} quintile, not the bottom one: a median runway of \result{QTwoRunway} months against
\result{QOneRunway} in the bottom quintile, and \result{QTwoUnderOne} percent under one month
against \result{QOneUnderOne} percent. The gradient is monotonic from the second quintile upward,
reaching \result{QFiveRunway} months and \result{QFiveUnderOne} percent at the top.

The explanation is not mysterious and it matters for policy design. Households in the bottom
quintile hold little debt and few of the commitments that a runway is measured against; households
in the second quintile have taken on the commitments of the households above them without the
income or the assets to carry them. A relief instrument targeted at the bottom of the wage
distribution misses the households with the least room, which is a design point rather than a
rhetorical one.

\begin{table}[htbp]
\centering
\caption{Liquid runway of working households, by wage quintile. Measured.}
\label{tab:buffers}
\begin{tabular}{lrrr}
\toprule
Wage quintile & Mean wage (USD) & Median months of runway & Under one month (per cent) \\
\midrule
Bottom quintile & 13,595 & 1.262 & 47.685 \\
Second quintile & 35,542 & 0.839 & 55.133 \\
Middle quintile & 55,416 & 1.090 & 49.732 \\
Fourth quintile & 84,909 & 1.402 & 43.234 \\
Top quintile & 227,078 & 1.808 & 36.398 \\
\bottomrule
\end{tabular}

\end{table}

\subsection{The difference between exposure types is a pay effect}

Households whose earners are physically exposed to automation look more fragile than households
whose earners are cognitively exposed. That difference does not survive a control for pay. Across
the cells we can compute both ways, reweighting the cognitively exposed group to the pay
distribution of the physically exposed group, and the reverse, leaves between
\result{PayGapSurvivingLow} and \result{PayGapSurvivingHigh} percent of the raw gap standing. On
one index and one survey the contrast reverses sign once pay is controlled.

The reading is that exposure type is largely a proxy for pay in these data, and that a household's
vulnerability to displacement is better predicted by what it earns than by what kind of work
produces the earnings. The direction of this result was confirmed in every cell by the independent
rebuild; the magnitudes were not, and we report the direction as the finding and the magnitudes as
ours alone.

\subsection{Attrition-led automation moves the private loss to the young}

Automation need not work through layoffs. It can work through hiring that does not happen, in
which case the people who bear it are not incumbents but entrants. We ran the first round three
ways: incumbents only, non-hiring only, and a sourced mix of the two. Table~\ref{tab:attrition}
gives all nine cells.

Two things come out of it. The first is compositional: under the non-hiring case, student loan
losses rise by \result{AttritionStudentLow} to \result{AttritionStudentHigh} percent against the
sourced mix and mortgage losses fall by \result{AttritionMortgageLow} to
\result{AttritionMortgageHigh} percent, because entrants carry student debt and rarely carry
mortgages. Auto losses barely move. The private loss is the same size and lands on different people
and different lenders.

The second is that the federal budget column does not move at all. It is identical across the three
incidence cases, to the cent, in every exposure group. A wage that is never earned costs the same
in tax whether it was never offered to an entrant or taken away from an incumbent. This is the
clearest statement in the paper of why the fiscal exposure is the structural one: it is invariant
to the channel through which displacement arrives, while the credit exposure is not.

\begin{table}[htbp]
\centering
\caption{First round at ten percent of the wage bill under three incidence cases, billions of
dollars. The federal budget column is identical across cases within each exposure group. Measured.}
\label{tab:attrition}
\begin{tabular}{llrrrr}
\toprule
Exposure group & Who is displaced & Federal budget & Mortgage & Auto & Student \\
\midrule
Physically exposed & incumbents only & 98.67 & 24.17 & 3.06 & 2.58 \\
Physically exposed & non-hiring & 98.67 & 20.49 & 3.61 & 3.70 \\
Physically exposed & sourced mix & 98.67 & 25.25 & 3.30 & 2.67 \\
Cognitively exposed, first index & incumbents only & 89.96 & 44.91 & 2.62 & 5.50 \\
Cognitively exposed, first index & non-hiring & 89.96 & 30.69 & 3.16 & 7.55 \\
Cognitively exposed, first index & sourced mix & 89.96 & 45.45 & 2.75 & 5.59 \\
Cognitively exposed, second index & incumbents only & 112.37 & 46.24 & 2.95 & 5.15 \\
Cognitively exposed, second index & non-hiring & 112.37 & 34.74 & 2.85 & 8.77 \\
Cognitively exposed, second index & sourced mix & 112.37 & 46.79 & 3.07 & 5.30 \\
\bottomrule
\end{tabular}

\end{table}

One gap in this engine is worth naming rather than leaving for a reader to find. It measures
displaced incumbents and, in the non-hiring case, jobs not offered to entrants who are already in
the data. People who are never hired at all are structurally invisible to it. External indicators
for recent graduates have been moving in ways consistent with that channel operating, and we report
it as a gap in our own measurement rather than as a finding of ours.
